\section{Meromorphic nonvanishing on simple spaces}\label{sec:nonvanishing}

The simple case of the induction needs a divisor representing a multiple of
the canonical bundle; its coefficients need not be nonnegative.  The
following result supplies precisely this starting point.

\begin{theorem}\label{nv:meromorphic}
Let \(X\) be a smooth connected compact Kähler manifold.  Suppose that
\(a(X)=0\) and that no positive-dimensional proper compact analytic
subvariety passes through a very general point of \(X\).  Then
\(K_X^{\otimes m}\) has a nonzero meromorphic section for some integer
\(m>0\).
\end{theorem}

The proof compares two ways of measuring poles.  After normalizing a big
class on \(X\) to have volume one, its pole order at a very general point is
bounded.  We construct a projective bundle over \(X^2\) whose volume grows
linearly with a parameter \(q\).  Restriction to the exceptional divisor of
the diagonal in the square of this bundle then produces a high-rank
subsheaf of a symmetric cotangent power.  A second diagonal forces
substantial vanishing of its determinant; descending that determinant
produces a point pole of order comparable to
\(q^{1/(2\dim X+1)}\), contradicting the point bound.
The argument uses only meromorphic sections of line bundles; it does not
require nonconstant meromorphic functions on \(X\).

For the proof, suppose that the conclusion fails.  Write
\[
 n=\dim X,\qquad \mathscr L=K_X,\qquad L=c_1(\mathscr L)
       \in H^{1,1}_{\BC}(X,\R).
\]
We may assume \(n\ge2\): a compact complex curve is projective, and the
assertion for a point is immediate.  Throughout this Section, an inequality
between real \((1,1)\)-classes is in pseudo-effective order:
\(\alpha\le\beta\) means that \(\beta-\alpha\) is pseudo-effective.  For a
divisor \(D\), \(\{D\}\) denotes \(c_1(\OO(D))\).  Pullbacks of classes will
occasionally be suppressed when the map is evident.

\subsection{Line subsheaves and point poles}

We first record the consequence of simplicity that controls all the
determinant lines used below.

\begin{lemma}\label{nv:slope}
Under the contrary hypothesis above, \(H^1(X,\OO_X)=0\), the group
\(\Pic(X)\) is countable, and \(L\ge0\).  If \(\mathscr H\) is a holomorphic
line bundle and
\(\mathscr H\longrightarrow\Omega_X^{\otimes k}\) is nonzero, where
\(k\ge0\), then
\begin{equation}\label{nv:slope-X}
 c_1(\mathscr H)\le kL .
\end{equation}
There is a complement of a countable union of proper analytic subsets of
\(X\) with the following further property.  For every point \(x\) in this
complement, let
\[
 a:Y=\Bl_xX\longrightarrow X,\qquad F=a^{-1}(x).
\]
Every nonzero line map
\(\mathscr H_Y\longrightarrow (a^*\Omega_X)^{\otimes k}\) satisfies
\begin{equation}\label{nv:slope-blowup}
 c_1(\mathscr H_Y)\le k a^*L .
\end{equation}
\end{lemma}

\begin{proof}
If the Albanese map of \(X\) is nonconstant, simplicity makes its image
\(n\)-dimensional: a positive-dimensional general fiber of smaller
dimension would be a forbidden subvariety.  At a point where the Albanese
map has rank \(n\), some \(n\) invariant one-forms have nonzero wedge after
pullback.  This gives a nonzero holomorphic section of \(K_X\), contrary to
our assumption.  Thus \(H^1(X,\OO_X)=0\).  The exponential sequence embeds
\(\Pic(X)\) into the countable group \(H^2(X,\Z)\).

The manifold \(X\) is not uniruled, by simplicity.  Ou's
non-pseudo-effectivity criterion for the canonical class therefore gives
\(L\ge0\) \cite[Theorem~1.1]{Ou2025}.  We recall how the slope and foliation
results in the same paper give the more precise inequality in
Equation~\eqref{nv:slope-X}.  Let \(\gamma\) be any class in the full dual of the
pseudo-effective cone, and use the slope
\[
 \mu_\gamma(\mathscr E)
   =\frac{c_1(\mathscr E)\cdot\gamma}{\rk\mathscr E}
\]
for torsion-free sheaves.  If
\(\mu_{\gamma,\min}(\Omega_X)<0\), the first Harder--Narasimhan piece
\(\mathscr T\subset T_X\) has strictly positive slope.  It is saturated and
semistable.  The tensor slope inequality shows that the bracket
\(\bigwedge^2\mathscr T\to T_X/\mathscr T\) is zero: the minimum slope of
its source is at least \(2\mu_\gamma(\mathscr T)\), whereas the maximum
slope of its target is smaller than \(\mu_\gamma(\mathscr T)\).
Consequently \(\mathscr T\) is a foliation.  Its dual has strictly negative
maximum slope and is not pseudo-effective by
\cite[Proposition~4.5]{Ou2025}.  Ou's foliation theorem
\cite[Theorem~1.4]{Ou2025} makes this foliation the tangent foliation of a
meromorphic fibration with compact general leaf closures.  Its rank is strictly
between zero and \(n\): rank \(n\) would give
\(\mu_\gamma(T_X)=-L\cdot\gamma/n\le0\).  The general leaf closure is
therefore a positive-dimensional proper subvariety through a general
point, again contradicting simplicity.

We have proved \(\mu_{\gamma,\min}(\Omega_X)\ge0\).  In its
Harder--Narasimhan filtration all quotient slopes are now nonnegative, so
\(\mu_{\gamma,\max}(\Omega_X)\le L\cdot\gamma\).  The tensor slope
inequality \cite[Lemma~4.4]{Ou2025} gives
\[
 c_1(\mathscr H)\cdot\gamma
 \le \mu_{\gamma,\max}(\Omega_X^{\otimes k})
 \le kL\cdot\gamma .
\]
Separation by the dual cone proves Equation~\eqref{nv:slope-X}.  For \(k=0\) this is
also immediate from the effective zero divisor of a nonzero map
\(\mathscr H\to\OO_X\).

It remains to choose \(x\) uniformly.  For any vector bundle
\(\mathscr V\), independent global sections are generically pointwise
independent.  Indeed, choose a maximal pointwise independent subfamily on
a dense open set.  Expressing any other section in that family by minors
gives meromorphic function coefficients.  They are constant because
\(a(X)=0\), so maximality among a linearly independent family forces all
its members to be pointwise independent.  Thus evaluation on
\(H^0(X,\mathscr V)\) is injective away from a proper analytic subset
(unless that vector space is zero, in which case there is no restriction).
Apply this to
\(\mathscr V=\mathscr H^{-1}\otimes\Omega_X^{\otimes k}\) for every
\(\mathscr H\in\Pic(X)\) and \(k\ge0\).  There are countably many such
bundles.

For \(x\) outside the resulting exceptional set, write any line bundle on
\(Y=\Bl_xX\) uniquely as
\(a^*\mathscr H\otimes\OO_Y(mF)\), \(m\in\Z\).  A line map from this bundle
to \((a^*\Omega_X)^{\otimes k}\), restricted away from \(F\), extends over
\(x\) to a section of
\(\mathscr H^{-1}\otimes\Omega_X^{\otimes k}\) by Hartogs' theorem.  If
\(m>0\), this section vanishes at \(x\), since
\(a_*\OO_Y(-mF)=\mathcal I_x^m\).  The choice of \(x\) rules this out.
Thus \(m\le0\), and Equation~\eqref{nv:slope-X} yields
\[
 c_1(a^*\mathscr H\otimes\OO_Y(mF))
 \le k a^*L+m\{F\}\le k a^*L.
\]
If a target has a finite filtration whose graded pieces are direct sums
of the indicated cotangent tensors, take the first nonzero associated
graded component of the line map and then a nonzero direct-sum component.
The preceding argument applies with the tensor order of that component.
\end{proof}

We will use several standard facts about volumes of real \((1,1)\)-classes.
Our normalization is \(\vol(\alpha)=\int\alpha^e\) for a nef class on an
\(e\)-fold.  Volume is continuous, homogeneous, monotone in
pseudo-effective order, invariant under modification, and its \(e\)-th
root is concave on the big cone.  Analytic Fujita approximation computes
it by Kähler parts on smooth projective modifications.  Here a
\emph{smooth projective modification} means a projective modification
whose source is smooth; the sources used here are compact Kähler
manifolds obtained by resolving coherent analytic ideals.  For a smooth
irreducible divisor \(D\) and a big class \(\alpha\), write
\(\operatorname{vol}_{\,|D}(\alpha)\) for the numerical restricted volume.
It is zero when \(D\) is contained in the non-Kähler locus
\cite[Theorem~1.1]{Vu2023}; otherwise it is
the supremum of the masses on \(D\) of restrictions of Kähler currents
with analytic singularities that are not generically singular on \(D\)
\cite[Lemma~2.7]{CollinsTosatti2018}.  The divisorial derivative and its
continuity on the big cone are
\begin{equation}\label{nv:volume-derivative}
 \frac{\dd}{\dd u}\vol(\alpha-u\{D\})
   =-e\,\operatorname{vol}_{\,|D}(\alpha-u\{D\});
\end{equation}
see \cite[Theorem~1.1]{Vu2023}.  We apply this formula only on smooth
manifolds while the varying class is big.

Here is a useful precise form of the approximation in the restricted
volume formula.  Resolve the log-ideal singularities of a current
restricted to \(D\).  Its pullback is an effective real divisor plus a
positive residual current with locally bounded potentials.  The latter
dominates a positive multiple of the pulled-back Kähler form.  A current
with locally bounded potentials has zero Lelong numbers, so Demailly
regularization makes its class, after subtracting that multiple, nef
\cite[Theorem~1.1]{DemaillyRegularization1992}.
On a projective modification there is an effective exceptional divisor
whose negative is relatively ample.  Subtracting a sufficiently small
multiple of this divisor from the residual class therefore makes that
class Kähler; add the same multiple to the divisor part.  Round all
divisor coefficients slightly upwards to rational numbers.  Openness of
the Kähler cone preserves the Kähler property.  These changes can be
arbitrarily small in top intersections, which compute the original mass
by the bounded-potential product formula.  Thus we may approximate a
restricted mass by
\begin{equation}\label{nv:rational-fujita}
 h^*(\alpha|_D)=\beta+\{D'\},\qquad
 \beta\ \text{Kähler},\quad D'\ge0\ \text{a rational divisor}.
\end{equation}
Moreover, \(D'\) has at least the log-ideal orders of the original
restriction on every further resolution.  We will use this last
property to turn poles into vanishing conditions.  Only \(D'\) is made
rational; the class \(\beta\) and the horizontal part of \(\alpha\) remain
real.

\begin{lemma}[A point pole bound]\label{nv:pole-threshold}
Let \(W\) be a smooth compact Kähler manifold of dimension \(e\ge2\),
\(\alpha\) a big real \((1,1)\)-class, and \(z\in W\).  Suppose a smooth
projective modification \(\mu:W'\to W\), which is an isomorphism near
\(z\), admits a decomposition
\[
 \mu^*\alpha=K+\{D\},\qquad K\ \text{Kähler},\quad D\ge0,
\]
where \(D\) misses the point \(z'\) over \(z\).  If the ordinary analytic
Seshadri constant \(\epsilon(K,z')\) is at least \(\eta>0\), then, on the
blowup \(b:\widehat W=\Bl_zW\to W\) with exceptional divisor \(G\),
\begin{equation}\label{nv:pole-formula}
 \sup\{u\ge0:b^*\alpha-u\{G\}\ge0\}
 \le \frac{\eta}{2}
     +2^{e-1}\vol(\alpha)\eta^{-(e-1)} .
\end{equation}
\end{lemma}

\begin{proof}
Set \(u_0=\eta/2\).  Blowing up \(z'\), the class
\(K-u_0\{G'\}\) is Kähler.  The Fujita decomposition is unchanged near
\(G'\), so it supplies a Kähler current for
\(\alpha_{u_0}=b^*\alpha-u_0\{G\}\) that is smooth near \(G\).
Its restriction there has class \(u_0c_1(\OO_{\PP^{e-1}}(1))\).
The restricted volume is consequently \(u_0^{e-1}\): the current gives
this lower bound, and the volume of the restricted class gives the
opposite bound.

Let \(\tau\) denote the left side of Equation~\eqref{nv:pole-formula}.  The classes
\(\alpha_u=b^*\alpha-u\{G\}\) are big for \(0\le u<\tau\), since
\(\alpha_0\) is big and the pseudo-effective cone is convex.  The
concave function \(f(u)=\vol(\alpha_u)^{1/e}\) satisfies, by
Equation~\eqref{nv:volume-derivative},
\[
 f'(u_0)=-\frac{u_0^{e-1}}{f(u_0)^{e-1}}.
\]
Its tangent line at \(u_0\) must remain positive up to \(\tau\).
Therefore
\[
 \tau\le u_0+\frac{f(u_0)^e}{u_0^{e-1}}
 \le u_0+\frac{\vol(\alpha)}{u_0^{e-1}},
\]
which is Equation~\eqref{nv:pole-formula}.
\end{proof}

Fix a Kähler class \(\omega\) on \(X\).  The class \(L\) is not big:
a big holomorphic line bundle would make \(X\) Moishezon, contrary to
\(a(X)=0\).  For \(t>0\) define
\begin{equation}\label{nv:normalization}
 P=r(L+t\omega),\qquad
 r=\vol(L+t\omega)^{-1/n}.
\end{equation}
Then \(P\) is a big real class, \(\vol(P)=1\), \(L\le P/r\), and
\(r\to\infty\) as \(t\downarrow0\).  Choose a smooth Fujita model
\(\mu:Y_P\to X\) with Kähler part \(P'\) satisfying
\begin{equation}\label{nv:P-prime}
 \mu^*P=P'+\{D_P\},\qquad v:=\int_{Y_P}(P')^n>\frac12.
\end{equation}
At a very general point \(y\) of this model there is no
positive-dimensional proper subvariety: its image would be one through a
very general point of \(X\), and \(y\) avoids the exceptional locus.
The ordinary Seshadri formula for a Kähler class,
\[
 \epsilon(K,y)
  =\inf_{\substack{W\ni y\\ \dim W>0}}
       \left(\frac{\int_W K^{\dim W}}
       {\operatorname{mult}_y W}\right)^{1/\dim W},
\]
therefore gives \(\epsilon(P',y)=v^{1/n}\ge2^{-1/n}\).
This is the ordinary nef/Kähler formula
\cite[Theorem~2.8]{Tosatti2016}; see also
\cite[Equation~(1.5)]{CollinsTosatti2018} and the
Kähler cone criterion of \cite{DemaillyPaun2004}.  We use this ordinary
formula in both applications below.  Lemma~\ref{nv:pole-threshold}, with
\(\vol(P)=1\), now yields a constant \(C_n\) depending only on \(n\)
such that, for a very general \(x\),
\begin{equation}\label{nv:point-bound}
 \tau(P,x):=\sup\{u\ge0:a^*P-u\{F\}\ge0\}\le C_n,
 \qquad a:\Bl_xX\to X.
\end{equation}
Here and below positive constants denoted \(c_n,C_n\) may be decreased or
increased from one occurrence to the next.  They are independent of all
parameters and choices of currents and modifications.

\subsection{A projective bundle with controlled volume}

We now produce a class whose volume grows, while its point pole threshold
grows much more slowly.  On \(X^2\), let \(\mathscr L_i\) and \(L_i\)
denote the pullbacks of \(\mathscr L\) and \(L\) from the \(i\)-th factor.
Use the quotient convention and put
\[
 \pi:Z=\PP_{X^2}(\mathscr L_1\oplus\mathscr L_2)\longrightarrow X^2,
 \qquad \xi=c_1(\OO_Z(1)),\qquad d=\dim Z=2n+1 .
\]
Thus \(\pi_*\OO_Z(m)=\Sym^m(\mathscr L_1\oplus\mathscr L_2)\) for
\(m\ge0\).  The zero divisor \(A_i\) of
\(\pi^*\mathscr L_i\to\OO_Z(1)\) has class \(\xi-L_i\); it is the
section on which \(\xi\) restricts to \(L_{3-i}\).  In particular
\(\xi=\{A_i\}+L_i\ge0\).

\begin{lemma}\label{nv:subvarieties}
For \(X\) and \(Z\) above, assume that \(K_X^{\otimes m}\) has no nonzero
meromorphic section for every integer \(m>0\).
Through a very general point of \(X^2\), the only positive-dimensional
proper compact irreducible analytic subvarieties are the two factor
slices.  Through a very general point of \(Z\), the only
positive-dimensional compact irreducible analytic subvarieties are a
fiber of \(\pi\), the full inverse images of the two factor slices, and
\(Z\) itself.
\end{lemma}

\begin{proof}
Let \(C\subset X^2\) be irreducible through a pair whose coordinates are
very general in \(X\).  Each projection image is a point or all of \(X\),
by properness and simplicity.  If exactly one is a point, \(C\) is the
corresponding full slice.  If both projections are surjective, test their
fibers at a general point of \(C\) whose coordinates are still very
general.  A positive-dimensional fiber must be the full other factor.
Thus either \(C=X^2\), or \(\dim C=n\) and both projections are
generically finite.

We show that subvarieties of the latter kind cannot sweep \(X^2\).
The space of compact \(n\)-cycles on a compact Kähler manifold has
countably many irreducible components, each compact
\cite[Theorem~1.1]{LiebermanCycles}.  On a component containing integral cycles
generically finite over both factors, integrality holds on a dense open
set and the two projection degrees remain positive, as can also be
tested by intersection with pulled-back Kähler forms.  Suppose the
incidence over one such component dominates \(X^2\).  Resolve the
parameter space and the dominating incidence component, obtaining maps
\[
 h:\mathcal Y\longrightarrow\mathcal S,\qquad
 F_i:\mathcal Y\longrightarrow X
\]
between smooth spaces.  For general \(t\in\mathcal S\), the fiber
\(Y_t=h^{-1}(t)\) is smooth, compact, and bimeromorphic to the integral
cycle, hence irreducible.  Its maps \(f_i=F_i|_{Y_t}:Y_t\to X\) are
generically finite.  Choose \(t\) also so that \(Y_t\) meets the dense
open set where \(d(F_1,h)\) is invertible and \(d(F_1,F_2)\) has rank
\(2n\).  The first assertion follows from generic finiteness of \(f_1\),
and the second from the assumed dominance in \(X^2\).

Our immediate aim is to turn this dominance into a meromorphic frame of
\(f_2^*T_X\).  Its determinant would meromorphically trivialize
\(f_2^*K_X^{-1}\), and hence also \(f_2^*K_X\).
On \(Y_t\), the equal-rank bundle map
\[
 d(F_1,h)|_{Y_t}:T_{\mathcal Y}|_{Y_t}
    \longrightarrow f_1^*T_X\oplus
                    (\OO_{Y_t}\otimes T_t\mathcal S)
\]
has a meromorphic inverse, given by its adjugate and determinant.
For fixed \(v\in T_t\mathcal S\), apply this inverse to \((0,v)\)
and then apply \(dF_2\).  The result is a meromorphic section of
\(f_2^*T_X\).  At a general point where \(d(F_1,F_2)\) has rank
\(2n\), the resulting map \(T_t\mathcal S\to f_2^*T_X\) is
surjective: in local coordinates \((F_1,h)\), it is the derivative
of \(F_2\) in the parameter directions with \(F_1\) fixed.  Choose
\(n\) fixed vectors \(v_i\) whose values there are independent.
Their wedge is a nonzero meromorphic section of \(f_2^*K_X^{-1}\)
on the irreducible \(Y_t\); its inverse trivializes \(f_2^*K_X\)
meromorphically.

Factor the proper generically finite map \(f_2\) as
\[
 Y_t\longrightarrow \overline Y_t
       \xrightarrow{\nu}X
\]
by Stein factorization.  Here \(\overline Y_t\) is normal and irreducible, the first
map is a modification, and \(\nu\) is finite of degree \(N>0\).
Meromorphic sections descend across a modification of a normal
space, so the section descends to a nonzero meromorphic section of
\(\nu^*K_X\).  Lemma~\ref{bir:finite-norm} gives a nonzero meromorphic
section of \(K_X^{\otimes N}\).  This norm is defined in local finite
analytic fraction algebras, so it remains available even though the
global meromorphic function field of \(X\) is just \(\C\).  It
contradicts our assumption.

Consequently the incidence image of each such cycle component is a
proper compact analytic subset of \(X^2\).  There are only countably
many components, so their union misses a very general pair.  This
proves the assertion for \(X^2\).

Now let \(W\subset Z\) be irreducible through a very general point.
Its image is a point, a full factor slice, or \(X^2\).  If its generic
relative dimension is one, it is the full inverse image of that
image.  Otherwise it is a multisection of the restricted
\(\PP^1\)-bundle.  Over a slice or over \(X^2\), a multisection is a
divisor with line bundle
\(\OO(k)\otimes\pi^*\mathscr H\), \(k>0\).  Its homogeneous equation
has coefficients in
\[
 H^0\bigl(S,\mathscr H\otimes
                 \mathscr L_1^{\otimes i}
                 \otimes\mathscr L_2^{\otimes(k-i)}\bigr),
 \qquad 0\le i\le k,
\]
where \(S\) is the relevant slice or \(X^2\), and a fixed-factor line is
understood as constant.  A single nonzero monomial cuts out only the
axis sections and vertical divisors.  A multisection not on an axis
therefore has two nonzero coefficients.  Their quotient is a
meromorphic section of a nonzero power of
\(\mathscr L_1\otimes\mathscr L_2^{-1}\).  Over a slice this is, up to
inversion and a constant line, a positive power of \(K_X\).  Over
\(X^2\), restriction to a general factor slice gives the same
conclusion.  This is impossible.  The axes themselves miss a very
general point of \(Z\), proving the claim.
\end{proof}

Let \(q\) tend to infinity through positive integers.  For each sufficiently
large \(q\), choose \(t=t(q)\) in Equation~\eqref{nv:normalization} so that
\(r=r(q)\ge q^2\), and use the resulting normalized class \(P=P(q)\).
Define the real class and the scale
\begin{equation}\label{nv:M-definition}
 M=P_1+P_2+q\xi,\qquad A_q=q^{1/d}.
\end{equation}

\begin{lemma}\label{nv:two-slot-volume}
For these choices, \(M\) is big and
\begin{equation}\label{nv:M-bounds}
 c_nq\le\vol(M)\le C_nq,\qquad
 \tau(M,z):=\sup\{u\ge0:b_z^*M-u\{F_z\}\ge0\}\le C_nA_q
\end{equation}
at a very general point \(z\in Z\), where
\(b_z:\Bl_zZ\to Z\) has exceptional divisor \(F_z\).
\end{lemma}

\begin{proof}
Pull \(Z\) to the Fujita model \(Y_P^2\) from
Equation~\eqref{nv:P-prime}, and put \(H=P'_1+P'_2\) on this pulled-back bundle.
There is a further smooth projective modification
\(\nu:\widehat Z\to\PP_{Y_P^2}(\mu^*\mathscr L_1\oplus
\mu^*\mathscr L_2)\) and a Fujita decomposition of the pullback of \(M\)
whose Kähler part is exactly
\begin{equation}\label{nv:exact-K}
 K=\frac12\nu^*H+\Theta ,
\end{equation}
where \(\Theta\) is Kähler and has degree \(q\) on a general vertical
line.  The modification and the divisor part miss that line.

Here is the construction.  It suffices to decompose \(H/2+q\xi\) as
\(\Theta\) plus an effective divisor, clean on a whole general vertical
line.  Since \(\OO(1)\) is relatively ample, a small positive class of
the form \(\delta(\xi+cH)\) is Kähler for some \(c>0\).  Choose
\(\delta>0\) so small that \(\delta<q\) and the remaining horizontal
part of \(H/2\) is Kähler.  Represent the remaining \(\xi\) first as
\(\{A_1\}+L_1\) and then as \(\{A_2\}+L_2\).  Regularize the
pseudo-effective classes \(L_i\) with analytic singularities, paying
their arbitrarily small negative errors from that horizontal Kähler
part.  This gives two Kähler currents in \(H/2+q\xi\), each smooth off
its own axis and a proper horizontal analytic set.  The maximum of
their potentials is locally bounded along an entire general vertical
line, since the two axes are disjoint.  Analytic regularization
preserving a smaller Kähler lower bound
\cite[Theorem~2.1(ii)]{Boucksom2004} gives a Kähler current with
analytic singularities missing that line.  Resolve its singularities
and make the small Kähler adjustment described before
Lemma~\ref{nv:pole-threshold}.  The divisor still misses a general
vertical line, so the residual class \(\Theta\) has degree exactly
\(q\) there.  Adding the other half of \(H\) and the pullbacks of the
divisor parts \(D_P\) gives Equation~\eqref{nv:exact-K}.  Its two
summands are nef, with \(\Theta\) Kähler, so every mixed intersection
used in the following bounds is nonnegative.

At a very general point of \(\widehat Z\), Lemma~\ref{nv:subvarieties}
lists the possible positive-dimensional subvarieties: the vertical line,
the strict transforms of the two full bundles over slices, and
\(\widehat Z\).  They have multiplicity one there.  By
Equation~\eqref{nv:exact-K} and nonnegativity of mixed intersections of nef
classes, their top intersections are respectively bounded below by
\begin{align}
 K\cdot(\text{vertical line})&=q, \notag\\
 \int_{\text{slice}}K^{n+1}
  &\ge \frac{n+1}{2^n}\,qv,\label{nv:K-intersections}\\
 \int_{\widehat Z}K^d
  &\ge \frac{d}{2^{2n}}\binom{2n}{n}\,qv^2. \notag
\end{align}
For example, the middle line is the term containing one factor
\(\Theta\) and \(n\) factors from the varying \(P'\); pushforward of
\(\Theta\) along a general vertical line is \(q\).  The last line is the
term with one \(\Theta\) and \(2n\) horizontal factors.  These
intersections give \(\vol(M)\ge c_nq\).  The ordinary Seshadri formula
and \(q\ge1\) give
\(\epsilon(K,z')\ge c_nq^{1/d}=c_nA_q\) at a very general \(z'\):
each possible dimension is at most \(d\), and every numerator in that
formula is at least \(c_nq\).

For the upper bound, subtract the axis \(A_1\).  For \(0\le u\le q\) put
\[
 M_u=M-u\{A_1\}
     =P_1+P_2+uL_1+(q-u)\xi .
\]
The endpoint \(M_q\) is pulled back from \(X^2\) and has zero volume on
the \(d\)-fold \(Z\).  It is pseudo-effective, and \(M_0\) is big by
the preceding construction, so \(M_u\) is big for \(u<q\).  The
restriction of \(M_u\) to \(A_1\simeq X^2\) is
\[
 (P+uL)_1+(P+(q-u)L)_2.
\]
The restricted volume along \(A_1\) is at most the volume of this
restriction.  For big classes \(\alpha_i\) on manifolds of dimensions
\(e_i\),
\begin{equation}\label{nv:product-volume}
 \vol(\pr_1^*\alpha_1+\pr_2^*\alpha_2)
   =\binom{e_1+e_2}{e_1}\vol(\alpha_1)\vol(\alpha_2).
\end{equation}
One can see this directly from the non-pluripolar product formula
\cite{BEGZ2010}: the
envelope with minimal singularities of a sum on a product is the sum
of the two envelopes, by testing the defining inequality on successive
slices.  Its top product has only the indicated binomial term.  Since
\(L\le P/r\) and \(r\ge q^2\), monotonicity and
Equation~\eqref{nv:product-volume} bound the restriction volume by
\[
 \binom{2n}{n}(1+u/r)^n(1+(q-u)/r)^n\le C_n.
\]
Integrating Equation~\eqref{nv:volume-derivative} from \(0\) to \(q\) gives
\(\vol(M)\le C_nq\).  Finally apply
Lemma~\ref{nv:pole-threshold} with \(e=d\),
\(\eta=c_nA_q\), and \(\vol(M)\le C_nq\).
Both terms on the right of Equation~\eqref{nv:pole-formula} are at most
\(C_nA_q\), since \(A_q^d=q\).  This proves Equation~\eqref{nv:M-bounds}.
\end{proof}

\subsection{A direct-image estimate}

The next lemma bounds a Kähler volume upstairs in terms of a class on the
base.  Its determinant formula will let Lemma~\ref{nv:slope} control that
base class.  The base need not be projective, and the horizontal class is
allowed to be real.

\begin{lemma}\label{nv:direct-image}
Let \(f:U\to Z_0\) be a surjective projective morphism of smooth connected
compact Kähler manifolds, with
\(\dim Z_0=b_0\ge1\) and \(\dim U=b_0+e\), \(e\ge1\).  Suppose
\[
 \beta=f^*G+c_1(\Lambda)
\]
is a Kähler class, where \(G\in H^{1,1}_{\BC}(Z_0,\R)\) and
\(\Lambda\in\Pic(U)\otimes\Q\).  Put
\[
 w=\int_{U_z}\beta^e>0
\]
on a general fiber.  For sufficiently large divisible integers \(j\),
let \(\mathcal F_j=f_*\OO_U(j\Lambda)\) and \(R_j=\rk\mathcal F_j\).
Here \(j\Lambda\) is an actual line bundle and
\(c_1(\mathcal F_j)\) means \(c_1(\det\mathcal F_j)\).  Then
\begin{align}
 R_j&=\frac{w}{e!}j^e+O(j^{e-1}),\label{nv:GRR-rank}\\
 B_0:=G+\lim_j\frac{c_1(\mathcal F_j)}{jR_j}
 &=\frac{f_*\beta^{e+1}}{(e+1)w}\ge0,\label{nv:GRR-base}\\
 \int_U\beta^{b_0+e}
 &\le (e+1)^{b_0}w\,\vol(B_0).\label{nv:mixed-bound}
\end{align}
The limit in Equation~\eqref{nv:GRR-base} is a limit of real Bott--Chern classes.
\end{lemma}

\begin{proof}
The restriction of \(c_1(\Lambda)\) to each fiber is represented by the
restriction of the Kähler form \(\beta\).  The fiberwise criterion for
relative ampleness makes \(\Lambda\) relatively ample after clearing
denominators.  Relative Serre vanishing then kills the higher direct
images for all sufficiently large divisible \(j\).  Analytic
Grothendieck--Riemann--Roch \cite{LevyGRR}, in degrees zero and two, gives
\[
 R_j=\frac{f_*c_1(\Lambda)^e}{e!}j^e+O(j^{e-1}),\qquad
 c_1(\mathcal F_j)
 =\frac{f_*c_1(\Lambda)^{e+1}}{(e+1)!}j^{e+1}+O(j^e).
\]
The degree-zero pushforward is \(w\).  Expanding
\(\beta=f^*G+c_1(\Lambda)\) shows that
\[
 f_*\beta^{e+1}
 =f_*c_1(\Lambda)^{e+1}+(e+1)wG,
\]
because terms with at least two horizontal factors have negative
fiber degree after pushforward.  This proves the equality in
Equation~\eqref{nv:GRR-base}; the cohomological GRR equality is an equality in
Bott--Chern cohomology by the \(\partial\bar\partial\)-lemma on compact
Kähler manifolds.  Pushforward of the positive form
\(\beta^{e+1}\) is a positive closed \((1,1)\)-current, so \(B_0\) is
pseudo-effective.

To prove Equation~\eqref{nv:mixed-bound}, we need a nef class on the base.
The class \(B_0\) is presently only known to be pseudo-effective.  We
flatten \(f\) to obtain a nef class above the base and compare it
with the pullback of \(B_0\).  Take a smooth projective flattening
modification \(p:Z'_0\to Z_0\), the equidimensional main transform
\(\overline U\) of \(U\times_{Z_0}Z'_0\), and a resolution
\(U'\to\overline U\).  Write \(f':U'\to Z'_0\) and \(q:U'\to U\) for
the maps and \(\beta'=q^*\beta\).  The class
\[
 B'_0=\frac{f'_*(\beta')^{e+1}}{(e+1)w}
\]
is nef.  It is enough to show that the positive pushforward current
representing this class has zero Lelong numbers.  This current can
be computed by integration on the cycle \(\overline U\) of the smooth
form pulled back from \(U\).  At a base point, cover the compact fiber
by finitely many coordinate neighborhoods, each embedded in a product
of base coordinates and ambient coordinates.  In the mass over a base
ball of radius \(\delta\), after wedging with a base Euclidean form to
power \(b_0-1\), the integrand is bounded by a finite sum of projection
volume forms using \(b_0-1\) base coordinates and \(e+1\) generic
linear combinations of all coordinates of the ambient product, including
the remaining base direction.  These projections can be chosen finite on the
neighborhoods: fixing the \(b_0-1\) base coordinates leaves local
dimension at most \(e+1\), by equidimensionality, and generic ambient
coordinates finish a finite projection.  After shrinking to compact
subneighborhoods their degrees are bounded.  Change of variables
therefore bounds each integral by
\(O(\delta^{2(b_0-1)})\) times the measure of the remaining coordinate
range.  On each compact subneighborhood these ranges, taken over closed
base balls, are nested compact sets whose intersection is the image of
the central fiber.  That image has measure zero in the \(e+1\) coordinates,
because the fiber has dimension \(e\).  Continuity of finite measure from
above shows that the range measures tend to zero.  Thus the
mass is
\[
 o\bigl(\delta^{2(b_0-1)}\bigr).
\]
This also covers \(b_0=1\), when it asserts absence of an atom.  The
pushforward current has zero Lelong numbers at every point, and
Demailly regularization \cite[Theorem~1.1]{DemaillyRegularization1992}
proves that its class \(B'_0\) is nef.

Pushforward under \(p\) gives \(p_*B'_0=B_0\).  For a modification between
smooth compact Kähler manifolds,
the kernel of pushforward on real Bott--Chern \((1,1)\)-classes is generated
by the classes of its exceptional prime divisors.  Since
\(p_*p^*B_0=B_0\), it follows that \(B'_0-p^*B_0\) is an exceptional real
divisor class.  It is
\(p\)-nef because \(B'_0\) is nef.  Apply relative negativity to this
exceptional real divisor.  Locally over the base, the usual proof for a
projective modification cuts by general hyperplanes to a surface and
uses the negative definite intersection matrix of exceptional curves;
it forces every coefficient of a relatively nef exceptional divisor to
be nonpositive.  Thus
\[
 B'_0=p^*B_0-\{D_0\},\qquad D_0\ge0\ \text{\(p\)-exceptional}.
\]
In particular
\begin{equation}\label{nv:flattened-volume}
 \int_{Z'_0}(B'_0)^{b_0}
  =\vol(B'_0)\le\vol(p^*B_0)=\vol(B_0).
\end{equation}
Choose a Kähler class \(\omega'\) on \(Z'_0\) and put
\(C=B'_0+\varepsilon\omega'\).  For \(0\le k\le b_0\), set
\[
 I_k=\int_{U'}(\beta')^{e+k}(f'^*C)^{b_0-k}.
\]
These are mixed intersections of nef classes.  The first two satisfy
\[
 I_0=w\int_{Z'_0}C^{b_0},\qquad
 I_1=(e+1)w\int_{Z'_0}B'_0C^{b_0-1}
       \le(e+1)w\int_{Z'_0}C^{b_0}.
\]
The mixed nef inequalities make the sequence \(I_k\) log-concave.
Every \(I_k\) is positive: on the dense open where \(q\) is a local
biholomorphism and \(f'\) is a submersion, \(\beta'\) is positive
definite and \(f'^*C\) has rank \(b_0\), so the defining top form
is strictly positive.  The successive ratios are therefore at most
\(I_1/I_0\le e+1\).  Thus
\[
 \int_U\beta^{b_0+e}=I_{b_0}
   \le(e+1)^{b_0}w\int_{Z'_0}C^{b_0}.
\]
Letting \(\varepsilon\downarrow0\) and using
Equation~\eqref{nv:flattened-volume} proves Equation~\eqref{nv:mixed-bound}.
\end{proof}

We will also use the determinant formula without its volume bound.

\begin{corollary}\label{nv:slope-descent}
Let \(Y\) be a smooth compact Kähler manifold and \(\mathscr E\) a
holomorphic vector bundle of rank at least two.  Write
\(\pi_{\mathscr E}:\PP_Y(\mathscr E)\to Y\) and
\(\zeta_{\mathscr E}=c_1(\OO_{\PP(\mathscr E)}(1))\).
Suppose that a real class \(D\) on \(Y\) satisfies
\[
 c_1(\mathscr H)\le kD
 \quad\text{whenever}\quad
 0\ne\bigl(\mathscr H\longrightarrow\mathscr E^{\otimes k}\bigr)
\]
for a line bundle \(\mathscr H\) and integer \(k\ge0\).
For a real class \(A\) on \(Y\) and a real number \(b\ge0\),
\begin{equation}\label{nv:slope-descent-equation}
 \pi_{\mathscr E}^*A+b\zeta_{\mathscr E}\ge0
 \quad\Longrightarrow\quad A+bD\ge0 .
\end{equation}
\end{corollary}

\begin{proof}
Choose a real class \(H_0\) on \(Y\) so that
\(\zeta_{\mathscr E}+\pi_{\mathscr E}^*H_0\) is Kähler; relative
ampleness of \(\OO(1)\) permits such a choice.  Take numbers
\(\delta\downarrow0\) with \(b+\delta>0\) rational.  Adding
\(\delta(\zeta_{\mathscr E}+\pi_{\mathscr E}^*H_0)\) to the
pseudo-effective class in Equation~\eqref{nv:slope-descent-equation} makes it
big.  A Fujita decomposition on a smooth projective modification
\(h:U\to\PP_Y(\mathscr E)\), with the divisor coefficients rounded
upwards as in Equation~\eqref{nv:rational-fujita}, has the form
\[
 \beta=f^*(A+\delta H_0)+c_1(\Lambda),\qquad
 \Lambda=(b+\delta)h^*\OO(1)-\OO_U(D')
       \ \text{in }\Pic(U)\otimes\Q,
\]
where \(f=\pi_{\mathscr E}h\), \(\beta\) is Kähler, and \(D'\ge0\)
is rational.  For divisible \(j\),
\[
 f_*\OO_U(j\Lambda)\ \subseteq\
           \Sym^{j(b+\delta)}\mathscr E .
\]
Its determinant of rank \(R_j\) consequently maps into
\(\mathscr E^{\otimes j(b+\delta)R_j}\).  The assumed line inequality
gives
\[
 \frac{c_1(\mathcal F_j)}{jR_j}\le(b+\delta)D.
\]
Lemma~\ref{nv:direct-image} makes
\(A+\delta H_0+\lim c_1(\mathcal F_j)/(jR_j)\) pseudo-effective.
Adding the preceding pseudo-effective difference shows that
\(A+\delta H_0+(b+\delta)D\) is pseudo-effective.  Let
\(\delta\downarrow0\).  This proves the corollary.  In this argument
only \(b+\delta\) and the coefficients of \(D'\) are rational; \(A\),
\(H_0\), and \(D\) remain real classes.
\end{proof}

\subsection{Restriction to the diagonal}

We now use the volume of \(M\) to obtain a high-rank subsheaf of a
symmetric cotangent power on \(Z\).  Distinguish the two copies of \(Z\)
by bracketed indices and
blow up their diagonal:
\[
 b:B=\Bl_{\Delta_Z}(Z\times Z)\longrightarrow Z\times Z,\qquad
 E=b^{-1}(\Delta_Z).
\]
Thus \(E=\PP_Z(\Omega_Z)\); its points are normal lines in \(T_Z\).
The dimensions are \(\dim B=2d\) and \(\dim E=2d-1\), and
\(E\to Z\) has relative dimension \(d-1\).
Put \(\zeta=c_1(\OO_E(1))\), so that \(\{E\}|_E=-\zeta\).
For \(s\ge0\) define
\[
 C_s=b^*(M_{[1]}+M_{[2]})-s\{E\},\qquad
 s_{\max}=\sup\{s\ge0:C_s\ge0\}.
\]
The class \(C_0\) is big, so \(s_{\max}>0\) and \(C_s\) is big for
\(0\le s<s_{\max}\).  Restricting a Kähler current with analytic
singularities to the fiber of the first projection at a very general
\(z\in Z\) gives the class \(b_z^*M-s\{F_z\}\).  The current can be
restricted for a general such \(z\), and Lemma~\ref{nv:two-slot-volume}
therefore gives
\begin{equation}\label{nv:smax}
 s_{\max}\le C_nA_q .
\end{equation}
Write \(v_E(s)=\operatorname{vol}_{\,|E}(C_s)\) for
\(0<s<s_{\max}\).  The restriction class is
\begin{equation}\label{nv:restriction-class}
 C_s|_E=2M+s\zeta .
\end{equation}

\begin{lemma}\label{nv:diagonal-mass}
There are constants \(c_n,C_n>0\) such that, for every sufficiently
large \(q\), one can choose a rational number
\[
 c_nA_q\le s\le C_nA_q
\]
and a Kähler current \(T\) in \(C_s\) with analytic singularities,
not generically singular on \(E\), with the following property.  On a
smooth projective modification \(h:U\to E\), its restricted mass has
a rational-divisor approximation
\begin{equation}\label{nv:selected-decomposition}
 h^*(2M+s\zeta)=\beta+\{D'\},
 \qquad \beta\ \text{Kähler},\quad D'\ge0\ \text{rational},
\end{equation}
which retains all log-ideal orders of \(T|_E\).  If
\(f:U\to Z\) is the natural map and
\[
 w=\int_{U_z}\beta^{d-1}
\]
on a general fiber, then
\begin{equation}\label{nv:selected-w}
 w\ge c_ns^{d-1}.
\end{equation}
\end{lemma}

\begin{proof}
First let \(s\) be any rational number in \((0,s_{\max})\) with
\(v_E(s)>0\), and use the approximation in
Equation~\eqref{nv:rational-fujita} for any current used
to compute this restricted volume.  In
Lemma~\ref{nv:direct-image}, the data for Equation~\eqref{nv:restriction-class}
are
\[
 b_0=d,\quad e=d-1,\quad G=2M,\quad
 \Lambda=s h^*\OO_E(1)-\OO_U(D')
           \ \text{in }\Pic(U)\otimes\Q.
\]
For large divisible \(j\), the associated sheaves satisfy
\begin{equation}\label{nv:symmetric-inclusion}
 \mathcal F_j=f_*\OO_U(j\Lambda)
       \subseteq\Sym^{js}\Omega_Z,\qquad
 R_j=\rk\mathcal F_j,\qquad 0<w\le s^{d-1}.
\end{equation}
The inclusion follows by pushing
\(\OO_U(-jD')\subseteq\OO_U\) through \(h\).
For the last inequality restrict the decomposition to a general
\(\PP^{d-1}\)-fiber: monotonicity of volume bounds the Kähler volume
there by \(\vol(s\zeta)=s^{d-1}\).

We claim that the class \(B_0\) of Equation~\eqref{nv:GRR-base} satisfies
\begin{equation}\label{nv:B0-bound}
 0\le B_0\le C_nM .
\end{equation}
Here is the determinant calculation, including its vertical sign.
The equality \(H^1(X,\OO_X)=0\) and the Künneth decomposition give
\(\Pic(X^2)=\pr_1^*\Pic(X)\oplus\pr_2^*\Pic(X)\).  Hence, for some
lines \(\mathscr Q_{j,i}\) on \(X\) and integer \(\ell_j\),
\[
 \det\mathcal F_j
  =\pi^*(\mathscr Q_{j,1}\boxtimes\mathscr Q_{j,2})
       \otimes\OO_Z(\ell_j),\qquad
 c_1(\det\mathcal F_j)=Q_j+\ell_j\xi,
\]
where \(Q_j=c_1(\mathscr Q_{j,1})_1+c_1(\mathscr Q_{j,2})_2\).
Taking the determinant of the inclusion in
Equation~\eqref{nv:symmetric-inclusion} gives a nonzero line
map into \(\Omega_Z^{\otimes k}\), \(k=jsR_j\).  It is first defined
where \(\mathcal F_j\) is locally free and extends over the
codimension-two complement.  Filter this tensor using
\[
 0\longrightarrow\pi^*\Omega_{X^2}\longrightarrow\Omega_Z
   \longrightarrow
   \OO_Z(-2)\otimes\pi^*(\mathscr L_1\otimes\mathscr L_2)
   \longrightarrow0 .
\]
A nonzero associated graded component has \(i\) relative factors and
\(k_1,k_2\) cotangent factors from the first and second factors of \(X^2\),
respectively, where
\(0\le i\le k\) and \(k_1+k_2=k-i\).  Its relative degree is \(-2i\).
Pushing the component to \(X^2\) forces
\[
 m_j:=-2i-\ell_j\ge0,\qquad \ell_j\le-2i\le0.
\]
A nonzero homogeneous monomial in
\(\Sym^{m_j}(\mathscr L_1\oplus\mathscr L_2)\), with exponent
\(0\le m_{j,1}\le m_j\) in the first factor, gives, on the two general slices,
\[
 c_1(\mathscr Q_{j,1})\le(i+m_{j,1}+k_1)L,\qquad
 c_1(\mathscr Q_{j,2})\le(i+m_j-m_{j,1}+k_2)L
\]
by Lemma~\ref{nv:slope}.  Each coefficient on the right is at most
\(k-\ell_j\).  Since \(L\ge0\), we obtain
\begin{equation}\label{nv:determinant-bound-Z}
 \ell_j\le0,\qquad
 Q_j\le(jsR_j-\ell_j)(L_1+L_2).
\end{equation}

Lemma~\ref{nv:direct-image} supplies a limit of the entire determinant
class divided by \(jR_j\).  The projective-bundle decomposition of
Bott--Chern cohomology gives separate limits \(Q\) and \(\ell\) of
the two displayed components.  Since
\[
 B_0=2(P_1+P_2)+(2q+\ell)\xi+Q\ge0,
\]
testing on a general vertical line gives \(2q+\ell\ge0\).
Equation~\eqref{nv:determinant-bound-Z} gives \(\ell\le0\) and
\(Q\le(s-\ell)(L_1+L_2)\).  Use \(\xi\ge0\), \(-\ell\le2q\),
\(L\le P/r\), \(r\ge q^2\), and \(s\le C_nA_q\le C_nq\).
They give
\[
 B_0\le2(P_1+P_2)+2q\xi+(s+2q)(L_1+L_2)
       \le C_nM,
\]
which proves Equation~\eqref{nv:B0-bound}.

Volume monotonicity, Lemma~\ref{nv:two-slot-volume}, and
Equation~\eqref{nv:mixed-bound} now imply
\begin{equation}\label{nv:restriction-mass-bound}
 \int_U\beta^{2d-1}\le C_nwq,\qquad
 v_E(s)\le C_nq\,s^{d-1}.
\end{equation}
For the second inequality, approximate the mass of each current
arbitrarily closely by \(\beta\) and use \(w\le s^{d-1}\), then take
the supremum over currents.  This proves it for rational \(s\);
continuity of divisorial restricted volume extends it to every
\(s\in(0,s_{\max})\).

At \(s_{\max}\) the class is on the boundary of the pseudo-effective
cone, so its volume is zero.  The product formula in
Equation~\eqref{nv:product-volume}, modification invariance, and the derivative
formula in Equation~\eqref{nv:volume-derivative} yield
\begin{equation}\label{nv:total-diagonal-mass}
 2d\int_0^{s_{\max}}v_E(s)\,\dd s
  =\vol(C_0)
  =\binom{2d}{d}\vol(M)^2
  \ge c_nq^2 .
\end{equation}
Choose a fixed small \(\theta_n>0\).  The contribution of
\(0<s<\theta_nA_q\), by Equation~\eqref{nv:restriction-mass-bound}, is at
most \(C_n\theta_n^d q^2\).  Fix \(\theta_n\) small enough that this
is less than half the last lower bound.  The remaining interval has
length at most \(C_nA_q\), by Equation~\eqref{nv:smax}.  Continuity therefore
permits a rational \(s\in[\theta_nA_q,s_{\max})\) such that
\[
 v_E(s)\ge c_nq^2/A_q=c_nqA_q^{d-1}
             \ge c_nq s^{d-1}.
\]
Choose a current with restricted mass at least three quarters of
\(v_E(s)\), and then an approximation as in
Equation~\eqref{nv:selected-decomposition} with
\(\int_U\beta^{2d-1}\ge v_E(s)/2\).  The first inequality of
 Equation~\eqref{nv:restriction-mass-bound} gives
 \(w\ge c_ns^{d-1}\), as claimed.
\end{proof}

For the selected \(s,T,h,\beta,D'\) and \(w\), retain the natural map
\(f:U\to Z\) and put
\[
 \Lambda=s h^*\OO_E(1)-\OO_U(D'),\qquad
 \mathcal F_j=f_*\OO_U(j\Lambda),\qquad R_j=\rk\mathcal F_j,
\]
where \(j\) is sufficiently large and divisible.  Here \(\Lambda\) is a
rational line, and Equation~\eqref{nv:symmetric-inclusion} gives
\(\mathcal F_j\subseteq\Sym^{js}\Omega_Z\).  With \(q\) and these selected
data fixed, Equations~\eqref{nv:GRR-rank} and \eqref{nv:selected-w} give
\begin{equation}\label{nv:selected-rank-fraction}
 \rk\Sym^{js}\Omega_Z=\binom{js+d-1}{d-1},\qquad
 \lim_j\frac{R_j}{\rk\Sym^{js}\Omega_Z}
       =\frac{w}{s^{d-1}}\ge c_n .
\end{equation}
The limit is through divisible integers.  Thus the first diagonal has
produced a subsheaf occupying a fixed positive proportion of the symmetric
power for all sufficiently large divisible \(j\).  The remaining proof
carries this rank bound to a modification of \(Z\) and converts it into a
large common order of vanishing for the determinant map on that model.

\subsection{A pole along the incidence}

We next locate a large pole of the selected current \(T\) along a smooth
incidence submanifold of \(B\).  This will impose vanishing conditions on the
symmetric-power subsheaf in Equation~\eqref{nv:symmetric-inclusion}.

Let \(U_0=\pi^{-1}(\Delta_X)\subset Z\), where \(\Delta_X\) is the
diagonal in \(X^2\).  Since
\(\mathscr L_1|_{\Delta_X}=\mathscr L_2|_{\Delta_X}=K_X\), there is a
canonical identification
\[
 U_0=X\times\PP^1.
\]
For \(\lambda\in\PP^1\) write \(D_\lambda=X\times\{\lambda\}\subset Z\).
Inside \(Z\times Z\) consider
\[
 \mathcal V_0
  =\{((x,x,\lambda),(y,y,\lambda)):x,y\in X,\ \lambda\in\PP^1\}
  \simeq X^2\times\PP^1.
\]
It meets \(\Delta_Z\) in \(\Delta_X\times\PP^1\).  The strict transform
\(V\subset B\) is
\(\Bl_{\Delta_X\times\PP^1}\mathcal V_0\); it is smooth, of dimension
\(d=2n+1\).  Its first projection is a smooth family over \(U_0\).
Over \(z=(x,x,\lambda)\) its fiber is
\[
 Y=\Bl_xD_\lambda\simeq\Bl_xX,
\]
and its intersection with \(E\) in that fiber is the projective space
of lines in \(T_zD_\lambda\), inside the projective space of lines in
\(T_zZ\).  These assertions follow either from the blowup of the
section \(y=x\) in this family or from the coordinates used below.

Let \(\rho_V:\Bl_VB\to B\) have exceptional divisor \(G_V\).  Denote
by \(b_V\ge0\) the generic divisorial pole of \(\rho_V^*T\) along
\(G_V\).  Equivalently, it is the generic log-ideal order of \(T\)
along \(V\), with the coefficient of the logarithm included.

\begin{lemma}\label{nv:incidence-pole}
The pole just defined satisfies
\begin{equation}\label{nv:pole-transfer}
 b_V\ge s-C_n .
\end{equation}
\end{lemma}

\begin{proof}
Remove \(b_V[G_V]\) from \(\rho_V^*T\).  Its residual positive current
can be restricted to \(G_V\): for analytic singularities removal of
the generic divisorial pole leaves a potential that is not identically
\(-\infty\) on that divisor.  Choose \(z=(x,x,\lambda)\in U_0\)
general enough for all restrictions and for
Lemmas~\ref{nv:slope} and \ref{nv:pole-threshold}.  On the bundle
\[
 G_V|_Y=\PP_Y(N^*_{V/B}|_Y),\qquad Y=\Bl_xX,
\]
this restriction is a positive current in the class
\begin{equation}\label{nv:incidence-restriction}
 a^*(2P+qL)-s\{F\}+b_V\zeta_V ,
\end{equation}
where \(\zeta_V=c_1(\OO_{\PP(N^*_{V/B}|_Y)}(1))\).
Indeed \(M|_{D_\lambda}=2P+qL\), the first projection to \(Z\) is constant on
\(Y\), and \(E|_Y=F\).  Also \(G_V|_{G_V}=-\zeta_V\), explaining the
positive sign of the last term.

The conormal bundle in Equation~\eqref{nv:incidence-restriction} has exact
sequences
\begin{align}
 0\longrightarrow\OO_Y^{\oplus n}
 &\longrightarrow N^*_{V/B}|_Y
 \longrightarrow a^*N^*_{D_\lambda/Z}\otimes\OO_Y(F)
 \longrightarrow0,\label{nv:incidence-conormal}\\
 0\longrightarrow\Omega_X
 &\longrightarrow N^*_{D_\lambda/Z}
 \longrightarrow\OO_X\longrightarrow0.\label{nv:small-conormal}
\end{align}
The first constant term is the conormal of \(U_0\) in the first copy of
\(Z\), evaluated at \(z\).  For the last term in
Equation~\eqref{nv:incidence-conormal}, the conormal of the strict transform of
\(D_\lambda\) in \(\Bl_zZ\) is
\(a^*N^*_{D_\lambda/Z}\otimes\OO_Y(F)\): in a blowup chart, a normal
coordinate is divided by an exceptional coordinate, giving precisely
the twist by \(F\) for conormals.  Equation~\eqref{nv:small-conormal}
is the conormal sequence for
\(D_\lambda\subset U_0\subset Z\); the diagonal in \(X^2\) has conormal
\(\Omega_X\), and the \(\lambda\)-normal direction is constant on
\(X\).

Filter a \(k\)-fold tensor of \(N^*_{V/B}|_Y\) by these sequences.  A
graded term has the form
\[
 \OO_Y(hF)\otimes(a^*\Omega_X)^{\otimes k'}
          \otimes(\text{a constant vector space}),
 \qquad 0\le k'\le h\le k.
\]
For any nonzero line map into that tensor, take a nonzero graded
component.  Lemma~\ref{nv:slope} gives
\[
 c_1(\mathscr H_Y)\le h\{F\}+k'a^*L
                 \le k(\{F\}+a^*L),
\]
since both \(\{F\}\) and \(a^*L\) are pseudo-effective.  Apply
Corollary~\ref{nv:slope-descent} to
Equation~\eqref{nv:incidence-restriction}, with
\(D=\{F\}+a^*L\).  We obtain
\begin{equation}\label{nv:incidence-base}
 a^*(2P+(q+b_V)L)-(s-b_V)\{F\}\ge0.
\end{equation}
If \(b_V\ge s\), the conclusion is immediate.  Otherwise
\(b_V<s\le C_nA_q\le C_nq\), and \(L\le P/r\) gives
\[
 \left(2+\frac{q+b_V}{r}\right)a^*P-(s-b_V)\{F\}\ge0.
\]
Use the point bound in Equation~\eqref{nv:point-bound}.  Since \(r\ge q^2\), its
consequence
\[
 s-b_V\le C_n\left(2+\frac{q+b_V}{r}\right)
\]
is bounded by a dimensional constant.  This proves
Equation~\eqref{nv:pole-transfer}.
\end{proof}

\subsection{Determinant vanishing and cancellation}

Blow up \(U_0\) in the base of \(E\):
\[
 p_+:Z^+=\Bl_{U_0}Z\longrightarrow Z,\qquad J_+=p_+^{-1}(U_0).
\]
This is the bundle \(Z\) pulled to \(\Bl_{\Delta_X}(X^2)\), since
\(U_0=\pi^{-1}(\Delta_X)\).  In particular \(Z^+\) is smooth.  Put
\[
 E^+=E\times_Z Z^+=\PP_{Z^+}(p_+^*\Omega_Z).
\]
Extend the incidence ideal from \(B\) to \(E\), and then to \(E^+\):
\[
 \mathcal J=(\mathcal I_V\cdot\OO_E)\cdot\OO_{E^+}.
\]
The products denote extension of ideals under the indicated maps.
Figure~\ref{nv:two-diagonals} locates \(V\cap E\) and this extension of
the incidence ideal to \(E^+\).

\begin{figure}[ht]
\centering
\begin{minipage}[c]{0.48\linewidth}
\centering
\small
\begin{tikzcd}[column sep=large,row sep=large]
 V \arrow[r,hook] \arrow[d,"\Bl_{\Delta_X\times\PP^1}"']
   & B \arrow[d,"\Bl_{\Delta_Z}"] \\
 \mathcal V_0 \arrow[r,hook] & Z\times Z
\end{tikzcd}
\[
 \begin{gathered}
 V_z=\Bl_xX,\\
 \begin{aligned}
 (V\cap E)_z&=\PP(T_z^*D_\lambda)\\
             &\subset E_z=\PP(T_z^*Z).
 \end{aligned}
 \end{gathered}
\]
\end{minipage}
\hfill
\begin{minipage}[c]{0.48\linewidth}
\centering
\small
\begin{tikzcd}[column sep=large,row sep=large]
 E^+ \arrow[r] \arrow[d] & E=\PP_Z(\Omega_Z) \arrow[d] \\
 Z^+ \arrow[r,"p_+=\Bl_{U_0}"] & Z
\end{tikzcd}
\[
 \begin{gathered}
 J_+=p_+^{-1}(U_0),\qquad U_0=X\times\PP^1,\\
 \mathcal J=(\mathcal I_V\cdot\OO_E)\cdot\OO_{E^+}.
 \end{gathered}
\]
\end{minipage}

\caption{Two geometric constructions used in the determinant estimate.
On the left, the strict transform \(V\subset B\) meets \(E_z\) in the
directions tangent to \(D_\lambda\).  Under the quotient convention,
the displayed projectivized cotangent spaces parametrize tangent lines.
On the right, \(E^+\) is the Cartesian base change
over \(Z^+=\Bl_{U_0}Z\), whose exceptional divisor is \(J_+\).
The restricted incidence ideal extends to \(\mathcal J\) on \(E^+\).
The ensuing local calculation identifies the normal parameter of \(J_+\)
among its generators.}
\label{nv:two-diagonals}
\end{figure}

The next lemma turns the pole bound in
Equation~\eqref{nv:pole-transfer} into a common vanishing order for the
coefficients of each determinant map.

\begin{lemma}\label{nv:determinant-order}
On a smooth projective modification \(h_+:U^+\to E^+\) one can choose
an approximation
\begin{equation}\label{nv:plus-decomposition}
 h_+^*(2p_+^*M+s\zeta)=\beta^++\{D^+\},
 \qquad \beta^+\ \text{Kähler},\quad D^+\ge0\ \text{rational},
\end{equation}
retaining the log-ideal orders of \(T|_E\), such that the general
fiber volume \(w^+\) of \(\beta^+\) over \(Z^+\) satisfies
\(w^+\ge w/2\).  Let \(f_+:U^+\to Z^+\) be the natural map and put
\[
 \Lambda^+=s h_+^*\OO_{E^+}(1)-\OO_{U^+}(D^+),\qquad
 \mathcal F_j^+=(f_+)_*\OO_{U^+}(j\Lambda^+),\qquad
 R_j^+=\rk\mathcal F_j^+ .
\]
For every sufficiently large divisible \(j\), write
\[
 \iota_j:\mathcal F_j^+\hookrightarrow p_+^*\Sym^{js}\Omega_Z,\qquad
 \varphi_j:\det\mathcal F_j^+\longrightarrow
       \bigwedge^{R_j^+}\!\bigl(p_+^*\Sym^{js}\Omega_Z\bigr)
\]
for the inclusion and its nonzero determinant map.  Here
\(\det\mathcal F_j^+=(\bigwedge^{R_j^+}\mathcal F_j^+)^{**}\), and the
determinant map extends across the complement of the locally free locus,
which has codimension at least two.
Localize a coordinate local ring of \(Z^+\) at the height-one prime of
\(J_+\), obtaining a discrete valuation ring \(R\) with uniformizer \(u\).
In local frames over \(R\), define
\[
 h_j=\min\{\ord_u(c):c\text{ is a nonzero coefficient of }\varphi_j\}.
\]
Equivalently, \(h_j\) is the minimum order of the maximal minors of
\(\iota_j\); it is independent of the frames.  If \(\sigma_{J_+}\) is the
canonical section of \(\OO_{Z^+}(J_+)\), then \(\varphi_j\) factors as
\[
 \det\mathcal F_j^+
 \xrightarrow{\ \cdot\sigma_{J_+}^{h_j}\ }
 \det\mathcal F_j^+(h_jJ_+)
 \xrightarrow{\ \widetilde\varphi_j\ }
 \bigwedge^{R_j^+}\!\bigl(p_+^*\Sym^{js}\Omega_Z\bigr),
\]
where \(\widetilde\varphi_j\) is holomorphic.  The common order satisfies
\begin{equation}\label{nv:determinant-order-equation}
 h_j\ge c_n\,jR_j^+s ,
\end{equation}
provided \(q\) is sufficiently large.
\end{lemma}

\begin{proof}
Blow up \(\mathcal J\) and take a common smooth projective resolution
with \(U\to E\) from Equation~\eqref{nv:selected-decomposition}.  Pull back
\(\beta\) and \(D'\).  Subtracting a sufficiently small rational
multiple of an effective exceptional divisor whose negative is
relatively ample makes the pulled-back Kähler class Kähler on this
resolution; add that multiple to the effective part.  Further upward
rational rounding may be
arbitrarily small.  This gives Equation~\eqref{nv:plus-decomposition} with all
original orders retained.  Its general fiber intersection is as close
to \(w\) as desired, so arrange \(w^+\ge w/2\).  Lemma~\ref{nv:direct-image}
and the effective divisor give
\begin{equation}\label{nv:plus-rank}
 \mathcal F_j^+\subseteq p_+^*\Sym^{js}\Omega_Z,\qquad
 R_j^+=\frac{w^+}{(d-1)!}j^{d-1}+O(j^{d-2}).
\end{equation}

We spell out the local order imposed on the polynomials in this
inclusion.  Near \(z=(x,x,\lambda)\), use coordinates
\((\mathbf x,\mathbf a,\lambda)\) on the first \(Z\), where \(\mathbf x\)
and \(\mathbf a\) each have \(n\) components and \(U_0=(\mathbf a=0)\).
Write the second coordinates as
\((\mathbf x+\mathbf h,\mathbf a+\mathbf k,\lambda+\mu)\), with
\(\mathbf h,\mathbf k\) each having \(n\)
components.  Then
\[
 \mathcal V_0=(\mathbf a=\mathbf k=\mu=0),\qquad
 \Delta_Z=(\mathbf h=\mathbf k=\mu=0).
\]
In a blowup chart meeting the lines tangent to \(D_\lambda\), take
\[
 h_1=v,\quad h_i=v\alpha_i\ (2\le i\le n),\quad
 k_i=vy_i\ (1\le i\le n),\quad \mu=vy_{n+1}.
\]
Here \(E=(v=0)\), while the strict transform is
\(V=(a_1=\cdots=a_n=y_1=\cdots=y_{n+1}=0)\).
In particular,
\begin{equation}\label{nv:ideal-before}
 \mathcal I_V|_E=(a_1,\ldots,a_n,y_1,\ldots,y_{n+1}).
\end{equation}
On a chart of the blowup of \(U_0\), write
\(a_1=u\) and \(a_i=u\tau_i\) for \(i\ge2\).  Along the general point
of \(J_+=(u=0)\), the ideal \(\mathcal J\) is exactly
\begin{equation}\label{nv:joint-ideal}
 (u,y_1,\ldots,y_{n+1}).
\end{equation}
Here \(u\) measures vanishing along \(J_+\), while the \(y_i\) measure
transverse fiber directions.  Thus a term of transverse degree \(\ell\)
can contribute at most \(\ell\) to the incidence order; the remaining
order must come from its coefficient in \(u\).  We now make this statement
precise.

Suppose locally that the singularities of \(T\) are
\(c\log|\mathfrak a|+O(1)\), with the logarithmic normalization for
which a divisorial coefficient is \(c\) times the ideal order.  If
\(k=\ord_V(\mathfrak a)\), then \(b_V=ck\).
The normal Taylor coefficients of degree less than \(k\) vanish on a
dense open set of \(V\), hence on \(V\) locally.  Thus
\(\mathfrak a\subseteq\mathcal I_V^k\) also near a general point of
\(V\cap E\).  Restricting to \(E\) and pulling to \(E^+\) gives the
corresponding order in Equation~\eqref{nv:joint-ideal}.  At its general center
the displayed generators are regular coordinates.  If \(H_{\mathcal J}\)
is the exceptional divisor of their blowup, its valuation is the
ordinary ideal-adic order:
\[
 \ord_{H_{\mathcal J}}(g)
  =\max\{k:g\in(u,y_1,\ldots,y_{n+1})^k\}.
\]
The common resolution dominates this blowup.  Pullback preserves the
effective difference between \(D'\) and the resolved log divisor,
and the later adjustments only increase the effective part.  Thus
\(D^+\) has coefficient at least \(b_V\) at this valuation.
Every polynomial image of a section of \(\mathcal F_j^+\) has
integral valuation at least \(\lceil jb_V\rceil\).  Locally bounded
remainders have zero order at this valuation.

There are \(n\) homogeneous coordinates along \(T_zD_\lambda\) and
\(n+1\) transverse homogeneous coordinates in \(T_zZ\).  Denote
these two groups by \(H_1,\ldots,H_n\) and
\(Y_1,\ldots,Y_{n+1}\).  For \(m=js\), a local polynomial in
\(\Sym^m(p_+^*\Omega_Z)\) has the form
\[
 \sum_{|\alpha|+|\gamma|=m}
          c_{\alpha,\gamma}(u)\,H^\alpha Y^\gamma.
\]
Work over the discrete valuation ring \(R\) used to define \(h_j\), choosing
the chart parameter \(u\) as uniformizer, and let \(\kappa\) be its residue
field.  The coefficients \(c_{\alpha,\gamma}\) lie in \(R\).
Equation~\eqref{nv:joint-ideal}
implies, monomial by monomial,
\begin{equation}\label{nv:coefficient-order}
 \ord_u(c_{\alpha,\gamma})
       \ge\max\{0,\lceil jb_V\rceil-|\gamma|\}.
\end{equation}
Indeed, on the chart \(H_1\ne0\), group the dehomogenized polynomial as
\[
 \sum_\gamma
 P_\gamma(H_2/H_1,\ldots,H_n/H_1)(Y/H_1)^\gamma,
 \qquad P_\gamma\in R[H_2/H_1,\ldots,H_n/H_1].
\]
Put \(b_j=\lceil jb_V\rceil\).  At the general center of
\((u,Y/H_1)\), order at least \(b_j\) forces
\(P_\gamma\) to be divisible by \(u^{b_j-|\gamma|}\) whenever
\(|\gamma|<b_j\).  Otherwise its first nonzero reduction would be a
nonzero polynomial over \(\kappa\), which stays nonzero in
\(\kappa(H_2/H_1,\ldots,H_n/H_1)\); in the associated graded ring
it would give a nonzero term of total \((u,Y/H_1)\)-degree less than
\(b_j\).  The transverse monomials there are independent.  Comparing
the internal polynomial coefficients over the base residue field
\(\kappa\) now gives Equation~\eqref{nv:coefficient-order}.

We compare the number of monomials with the rank estimate in
Equation~\eqref{nv:plus-rank}.  The
ambient rank and the number of monomials of transverse degree \(\ell\)
are
\[
 N_m=\binom{m+2n}{2n},\qquad
 N_{m,\ell}=\binom{\ell+n}{n}
             \binom{m-\ell+n-1}{n-1}.
\]
For fixed \(\delta\in(0,1)\), let \(H_m(\delta)\) count the monomials
with \(\ell>(1-\delta)m\).  Summing instead over their internal
degree gives
\[
 H_m(\delta)
 \le \binom{m+n}{n}\binom{\lceil\delta m\rceil+n}{n},\qquad
 \limsup_{m\to\infty}\frac{H_m(\delta)}{N_m}
       \le\binom{2n}{n}\delta^n.
\]
On the other hand, \(w^+\ge w/2\) and
Equation~\eqref{nv:plus-rank} show that the approximation on \(E^+\)
retains at least half the rank fraction in
Equation~\eqref{nv:selected-rank-fraction}:
\[
 \lim_{j\to\infty}\frac{R_j^+}{N_{js}}
       =\frac{w^+}{s^{d-1}}
       \ge\frac{w}{2s^{d-1}}\ge c_n>0.
\]
Fix \(\delta>0\), depending only on \(n\), so small that, for each fixed
\(q\) and its selected data, fewer than \(R_j^+/2\) monomials have transverse
degree greater than \((1-\delta)js\) for all sufficiently large divisible
\(j\).  For each
remaining monomial, Equation~\eqref{nv:coefficient-order} and
Lemma~\ref{nv:incidence-pole} give
\[
 \ord_u(c_{\alpha,\gamma})
   \ge j(s-C_n)-(1-\delta)js
   =j(\delta s-C_n)\ge\frac{\delta js}{2}.
\]
The last inequality holds once \(q\) is large, since
\(s\ge c_nA_q\to\infty\).

Over the same discrete valuation ring \(R\), the torsion-free sheaf
\(\mathcal F_j^+\) becomes a free module of rank \(R_j^+\).  Write its
inclusion into the symmetric power as a matrix with the monomials
as rows.  Every maximal minor uses at least \(R_j^+/2\) of the rows
whose coefficients have order at least \(\delta js/2\).  All maximal
minors therefore have order at least
\(\delta jR_j^+s/4\).  By the definition of \(h_j\), this proves
Equation~\eqref{nv:determinant-order-equation}.  Dividing the determinant
map by \(\sigma_{J_+}^{h_j}\) gives a holomorphic map at every point of
codimension one: the coefficients have the required order along \(J_+\),
and its local equation is a unit away from \(J_+\).  Hartogs' theorem
extends the divided map across the remaining set of codimension at least
two, giving the stated factorization through
\(\det\mathcal F_j^+(h_jJ_+)\).
\end{proof}

\begin{proof}[Completion of the proof of Theorem~\ref{nv:meromorphic}]
Fix \(q\) sufficiently large for the preceding lemmas, with its chosen
\(P,s,T\) and modifications.  Choose a very general point \(x\) in the first
factor of \(X^2\).  The slice of \(Z^+\) above \(x\) is
\[
 S=\PP_Y(\OO_Y\oplus a^*\mathscr L),\qquad
 a:Y=\Bl_xX\longrightarrow X,
\]
and \(J_+|_S\) is the pullback of \(F\).  We choose \(x\) so that
Equation~\eqref{nv:slope-blowup}, Equation~\eqref{nv:point-bound}, the restrictions of
the currents below, and all determinant maps for divisible \(j\)
can be tested on this slice.  These exclude at most countably many
proper analytic sets and the measure-zero exceptional sets for
current restrictions.

Write the restricted determinant line as
\[
 (\det\mathcal F_j^+)|_S
   =\pi_S^*\mathscr Q_j^+\otimes\OO_S(\ell_j^+),\qquad
 Q_j^+=c_1(\mathscr Q_j^+).
\]
Factoring the common zero of order \(h_j\) from the determinant
map gives a nonzero map whose source on \(S\) is
\[
 \pi_S^*(\mathscr Q_j^+\otimes\OO_Y(h_jF))
                 \otimes\OO_S(\ell_j^+).
\]
The plus sign of \(h_jF\) follows because division of the map by
the local equation of \(J_+^{h_j}\) enlarges its source from
\(\det\mathcal F_j^+\) to
\(\det\mathcal F_j^+\otimes\OO(h_jJ_+)\).
Its target embeds into the \(jsR_j^+\)-fold tensor of the
restriction of \(p_+^*\Omega_Z\).  In particular we use the
cotangent bundle of the original \(Z\), with its restricted
filtration
\[
 0\longrightarrow
 \pi_S^*(\OO_Y^{\oplus n}\oplus a^*\Omega_X)
 \longrightarrow (p_+^*\Omega_Z)|_S
 \longrightarrow
 \OO_S(-2)\otimes\pi_S^*a^*\mathscr L
 \longrightarrow0.
\]
The same homogeneous-monomial calculation as in
Equation~\eqref{nv:determinant-bound-Z}, now applying
Equation~\eqref{nv:slope-blowup} to the second factor of \(X^2\), gives
\begin{equation}\label{nv:determinant-bound-slice}
 \ell_j^+\le0,\qquad
 Q_j^++h_j\{F\}
       \le(jsR_j^+-\ell_j^+)a^*L.
\end{equation}
More explicitly, if the chosen graded term has \(i\) relative
factors, the pushed polynomial has degree
\(-2i-\ell_j^+\ge0\).  Its exponent in \(a^*\mathscr L\), together
with the \(i\) relative factors and the cotangent order from the second
factor of \(X^2\), is at most \(jsR_j^+-\ell_j^+\).  The cotangent
factors from the first factor of \(X^2\) are constant.  This proves the
displayed inequality using
the tensors of \(a^*\Omega_X\), without introducing \(\Omega_Y\)
or \(\Omega_{Z^+}\).

Apply Lemma~\ref{nv:direct-image} to the decomposition in
Equation~\eqref{nv:plus-decomposition}.  Equation~\eqref{nv:GRR-base}
shows that the limits of the restricted determinant components divided by
\(jR_j^+\) exist; denote them by \(Q^+\) and \(\ell^+\).  Divide
Equation~\eqref{nv:determinant-bound-slice} by \(jR_j^+\), use
Equation~\eqref{nv:determinant-order-equation}, and pass to the closed
pseudo-effective cone.  We obtain
\begin{equation}\label{nv:slice-upper}
 \ell^+\le0,\qquad
 Q^++c_ns\{F\}\le(s-\ell^+)a^*L.
\end{equation}
The pseudo-effective class in Equation~\eqref{nv:GRR-base}, restricted to the
very general slice \(S\), is
\begin{equation}\label{nv:slice-psef}
 \pi_S^*(Q^++2a^*P)+(\ell^++2q)\xi\ge0.
\end{equation}
Testing on a general vertical line gives \(\ell^++2q\ge0\).
Let \(A\subset S\) be the axis of
\(\PP_Y(\OO_Y\oplus a^*\mathscr L)\) with divisor class \(\{A\}=\xi\).
On \(A\simeq Y\), its normal class is \(\{A\}|_A=\xi|_A=a^*L\ge0\).
Add an arbitrarily small Kähler class to the class in
Equation~\eqref{nv:slice-psef}, and choose a Kähler current with analytic
singularities in the resulting big class.  Remove its generic divisorial
pole along \(A\) and restrict
the residual current to \(A\).  Adding back the removed nonnegative multiple
of \(a^*L\) preserves pseudo-effectivity.  Passing to the limit gives
\begin{equation}\label{nv:slice-lower}
 Q^++2a^*P+(\ell^++2q)a^*L\ge0.
\end{equation}

The slope bound in Equation~\eqref{nv:slice-upper} supplies the
pseudo-effective class \((s-\ell^+)a^*L-Q^+-c_ns\{F\}\).  Adding it to
Equation~\eqref{nv:slice-lower} cancels the entire normalized determinant
class restricted to the axis, \(Q^++\ell^+a^*L\), leaving
\[
 2a^*P+(s+2q)a^*L-c_ns\{F\}\ge0.
\]
Since \(L\le P/r\), the point bound in Equation~\eqref{nv:point-bound} implies
\begin{equation}\label{nv:contradiction}
 \frac{c_ns}{\,2+(s+2q)/r\,}\le C_n.
\end{equation}
The denominator is bounded independently of \(q\), because
\(s\le C_nq^{1/d}\) and \(r\ge q^2\).  The numerator tends to
infinity, because \(s\ge c_nq^{1/d}\).  Equation~\eqref{nv:contradiction}
is impossible for arbitrarily large \(q\).

For each \(q\), the choice of \(t(q)\), the rational \(s\), the current
\(T\), and the modifications precedes the limit in \(j\).  The monomial
cutoff \(\delta\) depends only on \(n\), while the required lower bound for
divisible \(j\) may depend on the fixed data.  We choose the very general
points after those data are fixed.
Thus every simultaneous general-point test above involves at most
countably many conditions, and no bound depends on the point.
The contradiction disproves the contrary hypothesis and proves
Theorem~\ref{nv:meromorphic}.
\end{proof}
